\section{병목 재진단: 페치가 아니라 데이터다}\label{sec:bottleneck}

\subsection{측정 인프라}

본 장과 V장의 모든 수치는 전용 벤치 하네스에서 나온
사이클 정밀 실측이다. 하네스는 SD카드$\cdot$JPEG$\cdot$LCD를 배제하고
플래시에 고정된 CIFAR-10 테스트 이미지 8장을 입력으로, 시드 고정
xorshift32 인코딩과 DWT\_CYCCNT(1사이클 해상도) 계측으로 동작한다.
워밍업 1장은 측정에서 제외한다. 각 측정은 당시의
\texttt{FLASH\_ACR}$\cdot$\texttt{SCB\_CCR} 레지스터 값과 ITCM 복사
바이트 수를 함께 기록하여, 빌드 라벨이 아니라 실리콘이 답한 값으로
측정 조건을 증명한다. 회귀 검증은 출력 스파이크
체크섬(\ref{sec:checksum}절)으로 한다.

선행 연구의 $2{,}407.0$ms는 1ms 해상도 시계와 표시 루프가 포함되고
입력 이미지 집합도 다른 조건의 수치이므로, 본 하네스로 동일 코드
경로를 재현한 기준선은 $2{,}328.5$ms이다(차이 3.3\%). 이후 모든
배속은 같은 하네스 안의 상대 비교로만 계산한다.

\subsection{절제 실측: 페치 최적화 2종은 0\%}

사전 하드웨어 조사는 두 가지 페치 최적화를 유망하게 지목했었다.
(i) 이 보드는 옵션 바이트에 의해 플래시의 ITCM 별칭(0x00200000)에서
부팅하는데 그 경로를 가속하는 ART 가속기가 꺼져
있었고(\texttt{FLASH\_ACR}$=$0x006), 별도 마이크로벤치마크에서 ART
활성화는 해당 경로를 1.4배 이상 가속했다. (ii) I-캐시(4KB)를 넘는
코드는 ITCM-RAM(16KB, 0 대기)에서 실행 시 최대 2.5배 빨랐다.

그러나 추론 워크로드에 대한 절제 실측(표~\ref{tab:ablation},
V장)에서 \textbf{둘 다 측정 노이즈 이내($\le$0.12\%)로 무효과}였다:
ART 활성화 $2{,}328.5\rightarrow2{,}328.9$ms, 핫 커널 전체(이 빌드
기준 $4{,}392$B, 베니어 포함)의 ITCM 배치
$\rightarrow2{,}331.2$ms. 원인은 두 가지다.

첫째, \textbf{부팅 주소와 실행 주소는 다르다.} 리셋 벡터는 ITCM 별칭에서
읽히지만, 링커가 코드를 AXIM 주소(0x08000000)에 배치했으므로 벡터에
적힌 주소부터 모든 페치는 AXIM 경로로 나간다. ART는 ITCM 경로
전용이라 AXIM 페치에 관여하지 못한다(별도 실측에서 AXIM 대역폭은 ART
on/off에 무반응). 즉 명령 페치는 이미 L1 I-캐시가 담당하고 있었다.

둘째, \textbf{핫 명령 집합이 캐시보다 훨씬 작다.} 실행 시간의
대부분은 SMLALD 내적 루프(수십$\sim$수백 바이트)에서 소요되며, 이는
4KB I-캐시에 상주한다. 사전 조사의 ``큰 커널 2.5배''는 6KB 코드를
균일 순회하는 합성 조건이었고, 실제 워크로드의 명령 지역성과는
달랐다.

\subsection{남는 용의자: 데이터 로드}

페치가 아니라면 데이터다. dense 커널의 CPL2는 2픽셀 타일마다 퓨전
가중치 전체($147{,}456$B)를 다시 읽으므로 가중치 트래픽만 타임스텝당
$32\times147{,}456\mathrm{B}\approx4.7$MB, 이미지당 141.6MB다. 이
스트림은 4KB D-캐시를 계속 밀어내며, 공급원이 플래시일 때는 라인필
지연이 더 크다. 실제로 가중치$\cdot$LIF 상태$\cdot$스크래치를
DTCM/SRAM으로 재배치한 절제 변형은 $2{,}178.7$ms(+7\%)를 얻었다 ---
유의미하지만, \textbf{트래픽의 총량을 줄이지 않는 배치 개선의 상한}이
이 정도임을 보여준다.

트래픽 총량을 줄이는 길은 \ref{sec:sparsity}절의 미활용 자원, 즉
``읽어 와서 0을 곱하는'' 연산 자체의 제거다. 다음 장이 그 커널이다.
